% Transcription — fonds Alexandre Grothendieck, shelfmark 46, batch 5 (pages 81-94).
% Established from the facsimile archives/batches/46/batch-05.pdf (cut from a
% copy of 46.pdf fetched through the project relay,
% https://grothendieck-relay.onrender.com/source/46.pdf; PDF page k =
% archivists' page 80+k, no cover sheet), per the skill transcribe-grothendieck.
% The folder is ninety-four pages in five batches; this is the fifth and last
% (fourteen pages).
%
% Pass: Opus 5.5 (claude-opus-5-5), 2026-09-26 — first pass, unchecked against the pages by a human.
%
% Pages skipped, and why: none as unrelated. Every leaf of 81-94 is in his
% hand. Page 93 is a cover (grey folder, inscribed top right « Contre-exemples
% à passage au quotient », the last two words doubtful); it gets no \page
% marker and its words become a section heading with a \note.
%
% Typescript: none. No further leaf of the « Définition du morphisme
% Verschiebung » typescript of batch 2 (pages 23-26) occurs in pages 81-94,
% and no other typed leaf either; the « ORMONT - St-DIÉ - VOSGES » showing
% through pages 84, 88-91 is a letterhead watermark in the paper, not text.
% No typescript decision to make in this batch.
%
% Pages summarised: none. Struck boxed blocks (pages 84, 85, 91, 92) and the
% passages cancelled by long diagonal strokes (foot of page 86, top of page
% 87, lower half of page 91, the English lines at the foot of page 94) are
% given, as far as they read, in \struck{} or in a \note.
%
% His pagination: none.
%
% What the batch is about. Pages 81-92 continue « Petits fourbis
% cohomologiques » from batch 4 (page 80 ends on the twisted Γ-set M_f):
% the invariants M_f^Γ (« éléments f-automorphes »), independence of M_f from
% the cocycle up to cohomology (u_a m = am); averaging N_f u = Σ f(γ)γ(u) and
% Poincaré's zeta-fuchsian series; Example: H^1(Γ, A^*) = (e) for a
% finite-dimensional algebra A over a field K on which Γ acts faithfully
% (polynomial P, independence of automorphisms for K infinite, Lang for K
% finite), and H^1(Γ, K^*) = 0 by linear independence. « Généralisation »:
% O semi-local noetherian, Γ finite acting faithfully and transitively on
% the maximal ideals, decomposition group Γ_d, inertia Γ_i, A an O-algebra:
% H^1(Γ, A^*) -> H^1(Γ_d, A_m^*) -> H^1(Γ_i, A_m^*) injective, via
% completion (Â = ⊕ Â_{m_i}), a density argument for the kernel of
% φ(u) = (u g(γ) - f(γ)γ(u)), dévissage by the exact sequences
% e -> N_n -> (A/m^{n+1}A)^* -> (A/m^nA)^* -> e with N_n ≃ m^nA/m^{n+1}A,
% the twist γ_f x = f(γ)γ.x f(γ)^{-1}, reduction to Γ_i = {e} and
% H^1(Γ, A^*) = {e} there; then, Γ acting trivially on O/m and on A/mA, a
% Sylow p-subgroup Λ: H^1(Γ, A^*) -> H^1(Γ, (A/mA)^*) × H^1(Λ, A^*) is
% injective, proved by the same dévissage and the bijectivity of
% multiplication by ν = [Γ:Λ] on A. Page 94, under the cover
% « Contre-exemples à passage au quotient »: G = GP(N)_S acting on P^N_S,
% S = Spec k; if G acts freely on an invariant open V and V/G exists, the
% sheaf descended from L = O(N+1) ≃ (Ω^N)^{-1} is ample, W = V/G is
% quasi-projective, whence V/G does not exist in the situation considered;
% the page ends with cancelled English lines (« to free creation ... from
% the heavy weight of ... »).
%
% Notation fixed once for the batch, following batch 4: the twisted action
% is \gamma_f (he writes γ with a small arrow above, f below); the maximal
% ideal, a plain m in his hand, is \mathfrak{m}; his underlined O is
% \mathcal{O}; completions carry \hat; the groups of units A^*, (A/mA)^*
% are set as he writes them, with the star sometimes over the A.
%
% Readings to check first: pages 83 and 86 (the connective prose), the
% margins of pages 86, 90 and 91, the struck block of page 84 and the foot
% of page 94.
% Apparatus: 195 \ill{}, 45 \uncertain{}.

\documentclass[11pt,a4paper]{article}
% Le préambule commun à toutes les transcriptions du fonds Grothendieck.
%
% Il tient en une idée : **l'incertitude se compose, elle ne se résout pas.**
% Une transcription qui lisse un mot douteux a détruit la seule chose pour
% laquelle on la faisait — un lecteur doit pouvoir savoir, à chaque endroit,
% ce qui était sur le feuillet et ce qui a été ajouté. Les six macros qui
% suivent sont donc l'appareil critique, et non de la décoration.
%
% Les mêmes commandes sont reconnues par `scripts/render.mjs`, qui produit la
% vue de lecture du site. Une macro ajoutée ici doit l'être aussi là-bas,
% faute de quoi le rendu échouera bruyamment — ce qui est voulu : un rendu
% silencieusement faux est pire qu'un rendu absent.

\ProvidesPackage{grothendieck}[2026/08/08 Transcriptions du fonds Grothendieck]

\usepackage{amsmath,amssymb,amsthm}
\usepackage{mathrsfs}
\usepackage{stmaryrd}
% Les diagrammes commutatifs sont partout dans ce fonds : c'est la langue
% même de la géométrie algébrique de Grothendieck, et une transcription qui
% les rendrait en prose ne transcrirait rien.
\usepackage{tikz-cd}
% Français par défaut — c'est la langue des feuillets — mais l'anglais est
% chargé aussi : la traduction et le résumé sont des documents anglais, et la
% typographie française (espaces avant les deux-points) y serait une faute.
% Ces documents basculent par \selectlanguage{english} en tête de corps.
\usepackage[english,french]{babel}
\usepackage{fontspec}
\usepackage{geometry}
\geometry{a4paper,margin=25mm,marginparwidth=28mm}
\usepackage{marginnote}
\usepackage{xcolor}
% ulem, not soul. `soul`'s \st and \ul are built on a character-by-character
% scan that XeLaTeX's font machinery breaks: the document fails with an
% undefined control sequence rather than degrading. ulem does the same two jobs
% and is Unicode-engine safe. `normalem` stops it hijacking \emph, which the
% transcriptions use for Grothendieck's own emphasis.
\usepackage[normalem]{ulem}
\usepackage{hyperref}
\hypersetup{colorlinks=true,linkcolor=black,urlcolor=black,pdfborder={0 0 0}}

% --- Métadonnées du lot -----------------------------------------------------
% Reprises telles quelles par le script de rendu, qui les affiche en tête de
% la vue de lecture. Elles fixent ce que le fichier prétend couvrir : sans
% elles, une transcription flotte sans point d'attache dans un fonds de seize
% mille feuillets.

\newcommand{\folder}[1]{\gdef\grothFolder{#1}}
\newcommand{\batch}[1]{\gdef\grothBatch{#1}}
\newcommand{\leaves}[2]{\gdef\grothFirst{#1}\gdef\grothLast{#2}}
\newcommand{\foldertitle}[1]{\gdef\grothFoldertitle{#1}}
\gdef\grothFolder{?}\gdef\grothBatch{?}\gdef\grothFirst{?}\gdef\grothLast{?}\gdef\grothFoldertitle{}

% --- Le résumé --------------------------------------------------------------
%
% Il ouvre la lecture modernisée, sous le titre « Résumé », et il est composé
% en petit. Ce n'est pas un ornement : c'est le seul passage du document qui
% ne soit pas des mathématiques, et un lecteur doit pouvoir le reconnaître
% avant de l'avoir lu — puis le sauter s'il connaît déjà le sujet, ou s'y
% arrêter s'il ne le connaît pas. Le filet qui le suit marque la reprise du
% corps.

\newenvironment{resume}
  {\small\setlength{\parskip}{0.45em}}
  {\par\medskip\hrule\medskip}

% --- Mots-clés ---------------------------------------------------------------
%
% En anglais, en clôture du résumé de la lecture modernisée : le vocabulaire
% moderne sous lequel chercher ce que les feuillets construisent. Le manifeste
% du site (`npm run manifest`) les extrait pour étiqueter la cote entière —
% cette ligne est la source unique des tags, il n'y a pas de fichier à côté.
\newcommand{\keywords}[1]{\par\smallskip\noindent
  {\footnotesize\textsc{Keywords} --- \textit{#1}}\par}

% --- L'appareil critique ----------------------------------------------------

% Le numéro de feuillet, en marge. C'est l'ancre qui permet de régler un
% désaccord sur une formule en regardant *un* feuillet plutôt qu'une chemise
% de six cents. Le site s'en sert aussi pour tourner les pages du fac-similé
% quand on fait défiler la transcription.
\newcommand{\leaf}[1]{%
  \marginnote{\footnotesize\color{gray!70}\textsf{#1}}%
  \label{leaf:#1}\ignorespaces}

% Illisible : on ne devine pas. Un mot inventé qui se lit comme les autres est
% le pire résultat possible.
\newcommand{\ill}{\textcolor{red!60!black}{[\,\ldots\,]}}

% Lecture douteuse : le mot est donné, et signalé comme incertain.
\newcommand{\uncertain}[1]{\uline{#1}}

% Ajout de l'éditeur — une lettre manquante, un mot sous-entendu.
\newcommand{\add}[1]{\textcolor{blue!50!black}{[#1]}}

% Biffé par Grothendieck lui-même. Ce qu'il a rayé fait partie du document :
% une rature dit souvent où la pensée a changé de direction.
\newcommand{\struck}[1]{\sout{#1}}

% Note du transcripteur, sur l'état matériel du feuillet.
\newcommand{\note}[1]{\par\smallskip{\small\color{gray!60!black}\textit{[#1]}}\par\smallskip}

% Note marginale de Grothendieck, distincte du corps du texte.
\newcommand{\marginal}[1]{\par\smallskip\noindent\hspace{1em}%
  {\small\color{gray!60!black}#1}\par\smallskip}

% --- Titre ------------------------------------------------------------------

\AtBeginDocument{%
  \begin{center}
    {\large\bfseries Fonds Alexandre Grothendieck --- cote n\textsuperscript{o}~\grothFolder}\\[2pt]
    {\itshape\grothFoldertitle}\\[4pt]
    {\small feuillets \grothFirst--\grothLast\ (lot \grothBatch)}\\[2pt]
    {\footnotesize\color{gray!60!black}Transcription établie d'après le fac-similé numérisé
     par l'Université de Montpellier.\\
     Les lectures incertaines sont soulignées, les illisibles notées [\,\ldots\,],
     les ajouts entre crochets.}
  \end{center}
  \vspace{1em}}

\endinput


\folder{46}
\batch{5}
\pages{81}{94}
\dating{[vers 1966-vers 1972]}
\watermark{Édition de démonstration}
\foldertitle{Schémas en groupes et fibrés principau[x] (Divers) : notes manuscrites (s.d.), tapuscrit (s.d.)}

\begin{document}

\page{81}

\note{suite de la page 80 (lot 4) : le $\Gamma$-ensemble tordu $M_f$}
Quel est l'ens.\ $M_f^{\Gamma}$ ? C'est l'ensemble des $m \in M_f = M$ tels
que
\[
  f(\gamma)\,\gamma m = m \quad \text{pour tt } \gamma \qquad \text{i.e.}
\]
\[
  \boxed{\gamma m = f(\gamma)^{-1} m \quad \text{pour tt } \gamma \in \Gamma}
\]
On les \uncertain{appelle} \uncertain{parfois} éléments $f$-automorphes du
$G$-$\Gamma$-ensemble $M$.

Si $M$ est muni \add{en plus} de structures diverses invariantes par $\Gamma$
et $G$, alors il en est de même de $M_f$ (qui conserve les structures qu'il y
avait sur $M$). Si $M$ était p.ex.\ un $G$-groupe abélien, \struck{\ill{}} un
$\Gamma$-groupe abélien, il en est de même de $M_f$ etc. Dans ce cas, on peut
s'intéresser, plus généralement, aux $H^i(\Gamma, M_f)$.

Soit $f'$ un autre cocycle \struck{\ill{}} de $\Gamma$ à valeurs dans $G$.
Supposons $f' \sim f$, i.e.\ $f'(\gamma) = a^{-1} f(\gamma)\,\gamma a$
($a \in G$). On a alors un isomorphisme $u_a : M_{f'} \to M_f$, en
\struck{faisant correspondre} \add{\uncertain{posant}}
\[
  u_a m = am
\]
\note{il écrit l'isomorphisme $M_f \to M_{f^*}$, l'astérisque étant
peut-être un prime ; le calcul qui suit est celui de $M_{f'} \to M_f$}
Il faut vérifier
\[
  u_a\, \gamma_{f'} m = \gamma_f\, u_a m
\]
\[
  \text{i.e.} \qquad a f'(\gamma)\,\gamma m = f(\gamma)\,\gamma(am)
\]
\[
  \text{i.e.} \qquad a a^{-1} f(\gamma)\,\gamma(a)\,\gamma(m) =
  f(\gamma)\,\gamma(am) \ ;
\]
c'est vrai.

Si on veut que $M_{f} \simeq M_{f'}$ pour tt $G$-$\Gamma$-ensemble $M$, il
\emph{faut} que $f \sim f'$ comme on voit en prenant $M = G$, $G$ y opérant par
translations à gauche.

\page{82}

Supposons que $M$ soit un groupe abélien.
La façon habituelle de construire des éléments de $M_f^{\Gamma}$ (du moins en
supposant $\Gamma$ \emph{fini}) est par le procédé des moyennes
\[
  N_f\, m = \sum_{\gamma \in \Gamma} \gamma_f\, m
  = \sum_{\gamma \in \Gamma} f(\gamma)\,\gamma.m
\]
Supposons en particulier que $G$ soit l'ens.\ des éléments inversibles d'un
\emph{anneau} $A$, et que $M$ soit un $A$-module. Soit
\struck{\ill{} \ill{} \ill{} $A$} \add{\ill{} \ill{} $u \in A$}
\struck{définissons}
\[
  N_f u = \sum_{\gamma \in \Gamma} f(\gamma)\,\gamma(u)
\]
Alors
\[
  x \in M_f^{\Gamma} \qquad \text{i.e.} \qquad f(\gamma)\,\gamma(x) = x
  \quad \text{si } x = N_f u
\]
\note{deux mots noircis devant $x \in M_f^{\Gamma}$ et devant le second $x$ ;
à droite, entouré : « $N_f(u) \in A_f^{\Gamma}$ », puis « En particulier, \ill{} »}
[Si le groupe $\Gamma$ est \uncertain{infini}, on \uncertain{essaie} de
\uncertain{faire} des moyennes infinies, quand on arrive à leur donner un sens.
C'est la \struck{\ill{}} construction des fonctions zeta-fuchsiennes de
Poincaré]

Si on peut s'arranger pour trouver $u$ tel que $x$ soit inversible, on aura
$f(\gamma) = x\,\gamma(x)^{-1}$ donc $f \sim 0$.

\emph{Exemple} \quad Supposons que $\Gamma$ \struck{\ill{}} opère
\emph{fidèlement} dans \add{un} corps $K$. Soit $A$ une algèbre de
dimension finie sur $K$, et supposons que $\Gamma$ opère dans $A$, de façon
que
\[
  \gamma(\lambda a) = (\gamma.\lambda)(\gamma.a)
\]
Je dis que alors $H^1(\Gamma, G) \simeq (e)$, où $G$ est le groupe des
éléments inversibles de $A$. D'après ce qui précède, \ill{} \uncertain{fini},
il suffit de trouver un $u \in A$ tel que
$\sum f(\gamma)\,\gamma u$ soit inversible, ce qui s'exprime par une équation

\page{83}

\[
  P\Bigl(\sum f(\gamma)\,\gamma a\Bigr) \neq 0
\]
où $P$ est un polynôme \uncertain{à coeff.}\ dans $K$ \struck{défini}
défini sur l'espace $K$-vectoriel $A$, par rapport aux coordonnées
\struck{\ill{}} du point courant, par rapport \ill{} \ill{} d'une base. On
montre d'abord que $A$ a une base formée d'éléments invariants par $\Gamma$
[c'est un lemme, qui n'utilise pas la structure \add{$K$-$\Gamma$-espace}
vectoriel \struck{\ill{}} de $A$, et n'utilise pas que la dimension soit
finie]. Si $a = (a_i)_{1 \leq i \leq n}$, \ill{} \ill{} \ill{}
\struck{$P\bigl(\sum f(\gamma)\,\gamma a\bigr) = Q\bigl(f_1(\gamma), \ldots,
f_n(\gamma)\bigr)$ ;}
\[
  P\Bigl(\sum f(\gamma)\,\gamma(a_\gamma)\Bigr) =
  R\Bigl(\bigl\lbrace \gamma a^1_\gamma, \ldots, \gamma a^n_\gamma
  \bigr\rbrace_{\gamma \in \Gamma}\Bigr)
\]
\note{la première ligne est biffée ; à droite : « $R$ un polynôme en
$n|\Gamma|$ variables »}
Supposons $K$ infini. \ill{} \uncertain{l'application} \ill{} \ill{}
$a^i_\gamma = a^i$ pour tt $\gamma$ (et tout $i$), étant données quelconques
$a^1, \ldots, a^n \in K$, alors d'après le th.\ d'indépendance algébrique des
automorphismes, alors $R$ est identiquement nul, en particulier on a
\struck{$R(\lbrace a\ldots$}
\[
  P\Bigl(\sum f(\gamma)\,\gamma(a_\gamma)\Bigr) = 0
\]
\uncertain{quel} que soit le choix des $a_\gamma$, d'où
(faisant $a_{\gamma_0} = 1$ et $a_{\gamma} = 0$ si $\gamma \neq \gamma_0$)
$P(f(\gamma_0)) = 0$, absurde.
\note{il écrit $a_{\gamma_1} = 0$ si $\gamma \neq \gamma_0$}
/ Dans le cas où $K$ est fini, il faut utiliser le résultat de Lang\ldots..

\page{84}

\struck{Généralisation} \emph{Cas particulier} \quad $H^1(\Gamma, K^{*}) = 0$.
\note{la formule est soulignée d'un trait ondulé}
Ici, il suffit de vérifier qu'on peut trouver $x \in K^{*}$ tel que
$\sum f(\gamma)\,\gamma x \neq 0$, et l'indépendance \emph{linéaire} des
automorphismes est suffisante.

\emph{Généralisation} \quad Soit $\mathcal{O}$ un anneau semi-local
noethérien, et $\Gamma$ un groupe fini opérant \emph{fidèlement} sur
$\mathcal{O}$, et par \uncertain{transitivité} sur les idéaux maximaux. Soit
$\Gamma_d$ le \struck{\ill{}} groupe de stabilité d'un idéal maximal
$\mathfrak{m}$ de $\mathcal{O}$, donc $\Gamma_d$ opère dans
$\mathcal{O}/\mathfrak{m} = k$. Soit $\Gamma_i$ le noyau de la représentation
de $\Gamma_d$ \struck{dans $\mathcal{O}/\mathfrak{m}$} \add{par}
automorphismes du corps $\mathcal{O}/\mathfrak{m}$. Soit enfin $A$ une
$\mathcal{O}$-algèbre, \struck{\ill{} de type fini} \struck{Alors}
\note{« \ill{} de type fini » est ajouté au-dessus de la ligne puis biffé
avec elle, et un « Alors » biffé suit}
$\Gamma$ opère \ill{} tel que $\gamma(\lambda a) = (\gamma\lambda)(\gamma a)$.
Alors les applications \struck{$\ldots$ l'application naturelle}
\[
  H^1(\Gamma, A^{*}) \longrightarrow H^1(\Gamma_d, A^{*}_{\mathfrak{m}})
  \longrightarrow H^1(\Gamma_i, A^{*}_{\mathfrak{m}})
\]
\note{après le dernier terme, quelques signes biffés d'un gribouillis}
sont \emph{injectives}. [N.B.\ \struck{\ill{} \ill{}} dans le cas le plus
important, \struck{$A$} $\Gamma_i$ opère trivialement sur $A/\mathfrak{m}A$,
donc $H^1(\Gamma_i, (A/\mathfrak{m}A)^{*})$ est formé des classes
d'homomorphismes \ill{} \ill{} \add{de $A/\mathfrak{m}A$} \ill{} de
$\Gamma_i$ dans $(A/\mathfrak{m}A)^{*}$]

Soient $f, g \in Z^1(\Gamma, A)$, supposons qu'elles aient même
\uncertain{image} dans \struck{$A_{\mathfrak{m}}$}
\add{$H^1(\Gamma_d$}
\struck{[i.e.\ il existe $a, b \in A$, $\lambda \in \complement\mathfrak{m}$,
$\mu \in \complement\mathfrak{m}$, tels que $\lambda^{-1}\mu^{-1} a b = 1$
dans $A_{\mathfrak{m}}$, i.e.\ il existe $\nu \in \complement\mathfrak{m}$
tel que $\nu ab = \nu\lambda\mu$, et tel que
$g(\sigma) = (a\lambda^{-1})^{-1} f(\sigma)\,\sigma(a\lambda^{-1})$ dans
$A_{\mathfrak{m}}$, i.e.\ \ldots]}
\note{le passage entre crochets est encadré et biffé d'un long trait
ondulé ; lecture des formules incertaine, les lettres $\lambda$, $\mu$, $\nu$
se distinguent mal, et $\complement\mathfrak{m}$ rend un signe qui peut
être « $\in \mathfrak{m}$ »}

\page{85}

Soient $\hat{\mathcal{O}}$, $\hat{A}$ les complétés de $\mathcal{O}$, $A$,
alors $\Gamma$ y opère de façon canonique. Or
\[
  \hat{\mathcal{O}} = \bigoplus_i \hat{\mathcal{O}}_{\mathfrak{m}_i}
  \qquad\qquad
  \hat{A} \simeq \bigoplus_i \underbrace{\hat{A}_{\mathfrak{m}_i}}_{B_i}
\]
et $\Gamma_d$ invarie $B_0$, tandis que $\Gamma$ mod $\Gamma_i$ permute les
$B_i$ \uncertain{entre} eux de façon simplement transitive.
\note{« $\Gamma$ mod $\Gamma_i$ » tel qu'écrit, là où l'on attend
$\Gamma_d$}
Je \uncertain{dis} qu'il résulte que
\[
  H^1(\Gamma, \hat{A}^{*}) \simeq H^1(\Gamma_d, \hat{A}_{\mathfrak{m}}^{*})
\]
d'après un lemme bien connu (valable même non commutatif). Il s'ensuit,
puisque $f$ et $g$ définissent la même \uncertain{classe} de
$H^1(\Gamma_d, \hat{A}_{\mathfrak{m}}^{*})$, \struck{qu'il existe un
$u \in A_{\mathfrak{m}}^{*}$} [en regardant le diagramme commutatif

\begin{tikzcd}[column sep=small, row sep=small]
H^1(\Gamma, A^{*}) \arrow[r] \arrow[d] & H^1(\Gamma_d, A_{\mathfrak{m}}^{*}) \arrow[d] \\
H^1(\Gamma, \hat{A}^{*}) \arrow[r, "\sim"] & H^1(\Gamma_d, \hat{A}_{\mathfrak{m}}^{*})
\end{tikzcd}

] que $f$ et $g$ ont même image dans $H^1(\Gamma, \hat{A})$ donc qu'il
existe un $u \in \hat{A}^{*}$ tel que l'on ait
\[
  g(\gamma) = u^{-1} f(\gamma)\,\gamma(u) \ :
\]
\struck{Par conséquent, on voit, \ill{} \ill{} \ill{} puissance
$\mathfrak{m}^N$ \ill{} \ill{} de $\mathcal{O}$, \ill{} \ill{} $v \in A$ \ill{}
$g(\gamma) \equiv v^{-1} f(\gamma)\,\gamma(v)$ mod $\mathfrak{m}^N A$}
\note{le passage est encadré et barré d'un long trait ondulé}

Considérons l'application \struck{\ill{}}
\[
  \varphi : A \longrightarrow A^{\Gamma}
\]
définie par
\[
  \varphi(u) = \bigl(u\,g(\gamma) - f(\gamma)\,\gamma(u)\bigr)_{\gamma \in \Gamma}
\]

\page{86}

Elle est \uncertain{linéaire} par rapport à $\mathcal{O}^{\Gamma}$. Or
$\mathcal{O}^{\Gamma}$ est un anneau semi-local noethérien, et $\mathcal{O}$
est \add{un module} de type fini sur $\mathcal{O}^{\Gamma}$ (bien
\uncertain{connu}\ldots), enfin la topologie naturelle sur $A$ (dans
$A^{\Gamma}$) est la même, qu'on considère $A$ comme
\add{$\mathcal{O}$-module} ou comme $\mathcal{O}^{\Gamma}$-module.
\note{il écrit « $A$ ($\text{dans } A^{\Gamma}$) », le $A^{\Gamma}$ étant le
but de $\varphi$}
Le noyau $N$ de l'application $\varphi$ a donc comme \uncertain{adhérence},
dans $\hat{A}$, le noyau $\hat{N}$ de l'application \uncertain{déduite} par
continuité. \ill{} \ill{} \ill{} \ill{} \ill{} \ill{} \ill{}. Soit
$v \in N$ tel que \emph{\uncertain{inversible}}. \struck{\ill{} \ill{} \ill{}}
$v \equiv u \ (\mathfrak{m}^r A)$ \ill{} \ill{} \ill{} \ill{} de
$\mathcal{O}$. Alors \ill{} est aussi inversible, et on a bien
\[
  g(\gamma) = v^{-1} f(\gamma)\,\gamma(v) .
\]
\note{les lettres $u$ et $v$ se distinguent mal sur cette page ; « $N$ » est
ajouté au-dessus de « noyau », « de type » au-dessus de « module »}
\marginal{\uncertain{digression} \ill{} \ill{} \ill{} si \ill{} $f \sim g$
mod $\mathfrak{m}^n A$ \struck{\ill{}} \ill{} $f \sim g$ \ill{}}
\note{dans la marge gauche, en quatre lignes obliques à la hauteur des
lignes 7 à 9 ; lecture très incertaine}

Cela prouve donc que
\[
  H^1(\Gamma, A^{*}) \to H^1(\Gamma_d, A_{\mathfrak{m}}^{*})
  \quad\text{et}\quad
  H^1(\Gamma, A^{*}) \to H^1(\Gamma, \hat{A}^{*})
  \quad\text{sont injectifs.}
\]
\note{les deux lignes sont réunies à gauche par un double trait vertical}

Supposons alors que $\mathcal{O}$ soit un anneau \emph{local} (donc
$\Gamma = \Gamma_d$), \struck{montrons que le diagramme} \ill{} que

\begin{tikzcd}[column sep=small, row sep=small]
H^1(\Gamma, A^{*}) \arrow[r] \arrow[d, "\varphi"'] & H^1(\Gamma_i, A^{*}) \arrow[d, "\varphi_i"] \\
H^1(\Gamma, A/\mathfrak{m}^{*}) \arrow[r, "\psi"] & H^1(\Gamma_i, A/\mathfrak{m}^{*})
\end{tikzcd}

\note{à gauche de la flèche $\varphi$, « \struck{est injectif} » biffé. Tout
ce bas de page, depuis « prouvons que », est barré de longs traits obliques}
\struck{Supposons que nous prouvions que $\varphi$ et $\psi$ sont
injectifs, alors il en résulte (\ill{} \ill{} $\Gamma_i$ \ill{}
$\varphi_i$ \ill{}}

\page{87}

\struck{aussi, et par suite aussi la première flèche \uncertain{horizontale}.}
\note{le haut de la page, jusqu'à « s'identifie au groupe additif
$\mathfrak{m}^n A / \mathfrak{m}^{n+1} A$ », est barré de longs traits
obliques, comme le bas de la page 86 ; on le transcrit néanmoins}

\struck{Supposons} Prouvons donc que $\varphi$ est injective. Prouvons
\uncertain{par} récurrence \uncertain{sur} $n \geq 1$ que
\[
  \text{si } f \sim g \ \mathrm{mod}\ \mathfrak{m}^n A
  \quad \text{alors} \quad f \sim g \ \mathrm{mod}\ \mathfrak{m}^{n+1} A
\]
Pour ceci, considérons la suite exacte de $\Gamma$-groupes
\[
  e \longrightarrow N_n \longrightarrow (A/\mathfrak{m}^{n+1}A)^{*}
  \longrightarrow (A/\mathfrak{m}^{n}A)^{*} \longrightarrow 0
\]
Si $x, y \in A^{*}$ \struck{\ill{}}, leurs images dans $A/\mathfrak{m}^n A$
sont égales si et seulement si $x - y \in \mathfrak{m}^{n}A$ i.e.\
$xy^{-1} - 1 \in \mathfrak{m}^{n}A$ \struck{\ill{} \ill{}} i.e.\
$xy^{-1} \in 1 + \mathfrak{m}^{n}A$, et donc \ill{} \ill{} l'identité
\struck{\ill{}} mod $\mathfrak{m}^{n+1}A$. Donc $N_n$ s'identifie au groupe
additif $\mathfrak{m}^n A / \mathfrak{m}^{n+1} A$ \ill{} \ill{} \ill{} que si
$f$, $g$ \ill{} tels que
\[
  g(\gamma) = u^{-1} f(\gamma)\,\gamma(u) \qquad \gamma \in \Gamma_i
  \qquad (*)
\]
alors $f \sim g$, i.e.\ \ill{} \ill{} d'abord faire opérer $\Gamma$
différemment sur $A$, \uncertain{par}
\[
  \gamma_f(x) = f(\gamma)\,\gamma.x\, f(\gamma)^{-1}
\]
\[
  \text{on a donc} \qquad
  \gamma_f(\lambda x) = f(\gamma)\,\gamma(\lambda)\,\gamma(x)\,f(\gamma)^{-1}
  = \gamma(\lambda)\,\gamma_f(x)
\]
donc $A_f$ est un $(\mathcal{O}, \Gamma)$-module. \struck{\ill{}} Cela étant,
l'égalité $(*)$ s'écrit
\[
  g(\gamma) f(\gamma)^{-1} = u^{-1} \gamma_f(u) \qquad \gamma \in \Gamma_i
\]
et on cherche $v$ tel qu'on ait
\[
  g(\gamma) f(\gamma)^{-1} = v^{-1} \gamma_f(v) \qquad
  \text{pour \emph{tt} } \gamma \in \Gamma .
\]

\page{88}

Or on sait que
\[
  h(\gamma) = g(\gamma) f(\gamma)^{-1}
\]
est un $1$-cocycle de $\Gamma$ à valeurs dans $A_f$.
On est ramené à prouver que s'il est \struck{trivial} \add{nul} sur
$\Gamma_i$, il est \struck{trivial} \add{nul} sur $\Gamma$. \ill{} \ill{}
\ill{} \ill{} \ill{} \ill{} \ill{}.
\note{« nul » est écrit deux fois au-dessus d'un mot biffé,
\uncertain{« trivial »}}
Cela nous ramène donc au cas où l'un des cocycles $f$, $g$ est l'unité,
soit $g$. Mais alors on utilise la suite exacte dans ce cas,
\uncertain{bien connue},
\[
  e \longrightarrow H^1(\Gamma/\Gamma_i, G^{\Gamma_i}) \longrightarrow
  H^1(\Gamma, G) \longrightarrow H^1(\Gamma_i, G)
\]
et on est ramené à prouver que
\[
  H^1(\Gamma/\Gamma_i, G^{\Gamma_i}) = \lbrace e \rbrace
\]
Or $G^{\Gamma_i} = (A^{\Gamma_i})^{*}$. Soit $\Gamma' = \Gamma/\Gamma_i$,
$A' = A^{\Gamma_i}$, \struck{Alors} $\mathcal{O}' = \mathcal{O}^{\Gamma_i}$.
On sait que $\mathcal{O}'$ est un anneau \add{local} noethérien, \ill{} \ill{}
\ill{} \ill{} \ill{} \ill{} isomorphe : \ill{} \ill{} que
$\Gamma/\Gamma_i = \Gamma'$ y opère \emph{fidèlement}. De plus, on sait que
$\mathcal{O}$ est un $\mathcal{O}'$-module \uncertain{noethérien}, d'où
résulte qu'il en est de même de $A$, donc de $A'$. $A'$ est une
$\mathcal{O}'$-algèbre, et $\Gamma'$ y opère de façon compatible avec ses
opérations sur $\mathcal{O}'$. On est ramené à prouver
\[
  H^1(\Gamma, A^{*}) = (e)
\]
\struck{\ill{}} dans le cas particulier où $\Gamma_i = \lbrace e \rbrace$,

\page{89}

i.e.\ $\Gamma$ opère \emph{fidèlement} sur le corps des restes. Or on sait déjà
(comme conséquence du th.\ d'indépendance algébrique des automorphismes,
resp.\ le lemme de Lang) que $H^1(\Gamma, A/\mathfrak{m}A^{*}) = (e)$.
\note{il écrit « \ldots\ automorphismes, resp! le \ldots\ » ;
$A/\mathfrak{m}A^{*}$, ici et plus bas, pour $(A/\mathfrak{m}A)^{*}$}
Prouvons alors par récurrence, sur $n \geq 1$ que
\[
  H^1(\Gamma, A/\mathfrak{m}^n A^{*}) = \lbrace e \rbrace
\]
Pour ceci, on \struck{suppose} le suppose fait pour $n \geq 1$, on écrit la
suite exacte
\[
  e \longrightarrow N_n \longrightarrow A/\mathfrak{m}^{n+1}A^{*}
  \longrightarrow A/\mathfrak{m}^{n}A^{*} \longrightarrow e
\]
et on constate que ces termes sont des $\Gamma$-groupes, $N_n$ est isomorphe à
$\mathfrak{m}^{n}A/\mathfrak{m}^{n+1}A$, donc à un espace vectoriel sur
$k = A/\mathfrak{m}$. On \ill{} \ill{} \ill{} (en tant que $\Gamma$-groupe)
isomorphe à $k^N$, donc $H^1(\Gamma, k^N) = H^1(\Gamma, k)^N = 0$ en vertu du
\emph{théorème} $H^1(\Gamma, k) = 0$ pour \uncertain{is.}
\note{« $k = A/\mathfrak{m}$ » tel qu'écrit, pour $\mathcal{O}/\mathfrak{m}$}
La suite exacte
\[
  H^1(\Gamma, N_n) \longrightarrow H^1(\Gamma, A/\mathfrak{m}^{n+1}A^{*})
  \longrightarrow H^1(\Gamma, A/\mathfrak{m}^{n}A)^{*}
\]
montre donc que l'on a
\[
  H^1(\Gamma, A/\mathfrak{m}^{n+1}A^{*})^{*} = \lbrace e \rbrace .
\]
\note{l'astérisque final est tel qu'écrit}

\page{90}

On, reproduisant un raisonnement déjà fait, on en conclut que
$H^1(\Gamma, A^{*}) = \lbrace e \rbrace$.

\uncertain{Cela} est \uncertain{clair} \uncertain{néanmoins} à étudier
\[
  H^1(\Gamma, A^{*})
\]
quand $\mathcal{O}$ est un \emph{anneau local}, $\Gamma$ opère
\emph{trivialement} sur $\mathcal{O}/\mathfrak{m}$, et $A$ est une
$(\mathcal{O}\text{-}\Gamma)$-algèbre. \emph{Supposons de plus que $\Gamma$
opère trivialement sur $A/\mathfrak{m}A$.}
\marginal{\ill{} \uncertain{plus} \uncertain{généralement}
\uncertain{supposons} que $A^{\Gamma} \to A/\mathfrak{m}A$ \ill{}
\emph{surjectif} [\ill{} \ill{} \ill{} \ill{} si $A = M_n(\mathcal{O})$,
$\Gamma$ \ill{} trivialement sur $\mathcal{O}/\mathfrak{m}$, \ill{}
\ill{} surjectif\ldots]}
\note{dans la marge gauche, en huit lignes obliques, reliées par un double
trait à « Supposons de plus que » ; « surjectif » est souligné deux fois}
Soit $p$ la caractéristique de $k = \mathcal{O}/\mathfrak{m}$, et soit
$\Lambda$ \struck{le} un sous-groupe de Sylow \struck{de type} $p$ de $\Gamma$
($\Lambda = \lbrace e \rbrace$ si la caractéristique est nulle).
Je dis que l'application
\[
  H^1(\Gamma, A^{*}) \longrightarrow H^1(\Gamma, A/\mathfrak{m}A^{*})
  \times H^1(\Lambda, A^{*})
\]
\struck{est bijective} est injective (en particulier, si
\ill{} \ill{} $\Lambda$ \ill{} \ill{} \ill{} $p$, \ill{}
\[
  H^1(\Gamma, A^{*}) \longrightarrow H^1(\Gamma, A/\mathfrak{m}A^{*})
  \quad\text{est injectif)}
\]
Soient $f$ et $g$ \struck{cocycles} deux cocycles dont les images dans
$H^1(\Gamma, A/\mathfrak{m}A^{*})$ et $H^1(\Lambda, A^{*})$ \ill{} \ill{}
\ill{} ; il faut montrer qu'ils sont $\sim$, \ill{} \ill{} \ill{} supposer
\add{\ill{}} $g(\gamma) = 1$ pour tt $\gamma$. \struck{Soit \ill{}} Ceci
revient alors \struck{\ill{}} \ill{} que $f$ est $\sim 0$ \ill{} \ill{}
$A/\mathfrak{m}^{n}A^{*}$ pour tout $n$. On raisonne par récurrence sur $n$,
supposons l'assertion prouvée pour un $n \geq 1$.
\note{le « $1$ » de $g(\gamma) = 1$ est surchargé ; « \ill{} » ajouté
au-dessus de la ligne, en fin de ligne}

\page{91}

La suite exacte de coh.\ montre que $f$ est dans l'image de
\[
  H^1(\Gamma, \mathfrak{m}^{n}A/\mathfrak{m}^{n+1}A) \longrightarrow
  H^1(\Gamma, A/\mathfrak{m}^{n+1}A^{*})
\]
\add{dans $A/\mathfrak{m}^{n+1}A^{*}$} i.e.\ est $\sim$ à un cocycle de la
forme \struck{\ill{}}
\[
  1 + g(\gamma)
\]
où
\[
  g(\gamma) \in \mathfrak{m}^{n}A \quad \text{satisfait à} \qquad
  g(\gamma\gamma') \equiv g(\gamma) + \gamma g(\gamma') \quad
  (\mathfrak{m}^{n+1}A)
\]
\struck{Par hypothèse, on} On peut dans le raisonnement supposer
$\mathfrak{m}^{n+1} = 0$. \struck{Donc} L'hypothèse \struck{\ill{}} signifie
que
\[
  1 + g(\gamma) = u^{-1}\gamma(u) \qquad \text{pour } \gamma \in \Lambda
\]
\[
  \text{i.e.} \qquad \gamma(u) - u = u\,g(\gamma) \qquad
  \text{pour } \gamma \in \Lambda .
\]
\struck{Or $\mathfrak{m}\,g(\gamma) = 0$, \ldots\
$u\,g(\gamma) = \varepsilon(u)\,g(\gamma)$ i.e.\
$g(\gamma) = \gamma(\varepsilon(u)^{-1}u) - \varepsilon(u)^{-1}u$
\ldots\ ce qui montre que $\Gamma$ opère trivialement sur
$\mathcal{O}/\mathfrak{m}$ !)}
\note{bloc encadré et barré de traits obliques ; le « $\mathfrak{m}$ » de
« $\mathfrak{m}\,g(\gamma) = 0$ » est peut-être un autre signe}
\struck{Cela signifie que \ill{} $g(\gamma)$ est $\sim 0$ dans
$H^1(\Lambda, A)$ (additif !). \ill{} Or, $A$ est un
$(\Gamma\text{-}\mathcal{O})$-module de longueur finie, \ill{} \ill{} \ill{}
\ill{} \ill{} \ill{} \ill{} $M$ \ill{} que l'application
$H^1(\Gamma, M) \to H^1(\Lambda, M)$ est \emph{injective}.}
\struck{\ill{} \ill{} $M$ \ill{} \ill{} \ill{} \ill{}}
\note{tout le bas de la page, depuis « Cela signifie », est barré de longs
traits obliques ; les deux dernières lignes sont en outre biffées chacune
d'un trait}
\marginal{\ill{} $\mathfrak{m}M = M$, \ill{} $\mathfrak{m}M = 0$
\ill{} $\Gamma$ \ill{} \ill{} $H^1(\Gamma, M) \to H^1(\Lambda, M)$
\ill{} \ill{} \ill{} \ill{} \ill{} \ill{} \uncertain{injectif}.}
\note{la marge gauche porte, verticalement et en partie barrée d'un trait
oblique, une note de huit lignes courtes ; lecture très incertaine}

\page{92}

\struck{Caractéristiques. Si \ill{} \ill{} $0 \to M' \to M \to M'' \to 0$
\ill{} \ill{} \ill{} \uncertain{surjectif} \ill{}}
\note{bloc encadré et biffé ; à droite, embrassé d'une accolade :
« provient des faits que, si »}
$\nu = [\Gamma : \Lambda]$, la multiplication par $\nu$ induit une
bijection dans $A$ (car ce \ill{} \ill{} \ill{} dans les
$\mathfrak{m}^{i}A/\mathfrak{m}^{i+1}A$).
\struck{Par suite, on \ill{} \ill{} \ill{} \ill{} un $v \in A$ \ill{}
$\gamma v - v = u\,g(\gamma)$ pour tout $\gamma \in \Gamma$}
\note{ces deux phrases, entre l'accolade et le bloc biffé, sont barrées de
traits obliques, la dernière encadrée et biffée}

Soit $v \in A^{\Gamma}$ tel que $\varepsilon(v) = \varepsilon(u)$, on a donc
\[
  g(\gamma) = \gamma(v^{-1}u) - v^{-1}u
\]
d'où $g \sim 0$ dans $H^1(\Lambda, A)$ (additif). Comme $\nu = [\Gamma :
\Lambda]$ est tel que $x \mapsto \nu x$ dans $A$ est \emph{bijectif} (car il
l'est dans les $\mathfrak{m}^{i}A/\mathfrak{m}^{i+1}A$, qui sont des espaces
vectoriels sur $k$) l'homomorphisme
\[
  H^1(\Gamma, A) \longrightarrow H^1(\Lambda, A)
\]
est \emph{injectif}, donc il existe $w \in A$ tel que
\[
  g(\gamma) = \gamma(w) - w \qquad \text{pour \emph{tout} } \gamma \in \Gamma .
\]
Donc \struck{$1 + g(\gamma) = \ldots$ $w g$} On \uncertain{peut} \ill{}
(\ill{} \ill{} \ill{} \ill{} $A^{\Gamma}$) supposer $w \in A^{*}$,
\struck{\ill{}} \uncertain{et} \ill{} \ill{} $\varepsilon(w) = 1$. Alors
\ill{}. \struck{Soit alors}
\[
  g(\gamma) = \varepsilon(w)\,g(\gamma) = w\,g(\gamma), \quad \text{d'où}
\]
\[
  g(\gamma) = w^{-1}\gamma(w) - 1 \quad \text{d'où}
\]
\[
  1 + g(\gamma) = w^{-1}\gamma(w)
\]
donc $1 + g(\gamma) \sim 0$ dans $H^1(\Gamma, A^{*})$.

\section*{Contre-exemples à passage au quotient}

\note{inscrit en haut à droite d'une chemise grise (page 93), qui ne reçoit
pas de numéro de page ; les mots « passage au quotient » sont d'une lecture
incertaine}

\page{94}

\marginal{Impossibilité des \uncertain{passages} \uncertain{aux quotients}}
\note{en haut à droite, dans un angle tracé à la plume}

$G = GP(N)_S$ opérant sur $\mathbb{P}^N_S$ \qquad $S = \operatorname{Spec} k$

$V$ ouvert de $\mathbb{P}^N_S$ invariant \uncertain{sur lequel} $G$ opère
librement, supposons que $V/G$ existe « effectivement ». Alors on montre que
le faisceau sur $V/G = W$ déduit de
$\underline{L} = \underline{\mathcal{O}}(N+1) \simeq
\bigl(\underline{\Omega}^{N}_{\mathbb{P}^N_S/S}\bigr)^{-1}$ par « descente »
est \emph{ample} sur $W$ [utiliser le fait que $G$ n'a pas d'homomorphismes
non \uncertain{triviaux} dans $\mathbb{G}_m$, et que tout ouvert affine de
$\mathbb{P}^N_S$ est complémentaire d'une hypersurface, \struck{\ill{}} donc
défini par une section d'un $\underline{L}^{\otimes m}$ ($m > 0$)]. A
fortiori, $W$ est \emph{quasi-projectif} ; \struck{\ill{}} donc si $G$ opère
\add{\uncertain{librement}} dans $V$ \ill{} \ill{} \ill{} \ill{} \ill{}
\ill{} ; $V/G$ n'existe pas.

\struck{\ill{} \add{so as} to free creation \ill{} from the \ill{} weight of
\ill{} concepts \ill{} and \ill{} \ill{} and \ill{} tools}
\note{ces lignes en anglais, au bas de la page, sont barrées de longs traits
obliques et portent plusieurs corrections interlinéaires, elles-mêmes
biffées ; on lit notamment « the heavy weight of », « conceptual
background and technical » (biffés) et « tools »}

\end{document}
